\subsection{The stable positional fifty-three factor and its coherent divisors}

\begin{statement}{E.123}{Proved}{Actual positional class norms, unit defects and a complete coherent-divisor image}
Keep the compatible roots of E.119--E.122, and put
\[
 \begin{gathered}
 K=\mathbb Q(\zeta_3),\quad \xi=\zeta_9,\quad
 \beta_d^{3^d}=53\zeta_3,\quad \beta_{d+1}^3=\beta_d,\\
 L_d=\mathbb Q(\beta_d),\quad B_d=L_d(\xi),\quad
 C_d=\operatorname{Cl}_S(L_d)\{3\},\quad
 V_d=\mathcal O_{L_d,S}^{\times},\qquad S\mid\{3,53\}.
 \end{gathered}
 \tag{E123.1}
\]
For every $d\ge2$, the horizontal extension $B_d/L_d$ is cyclic,
unramified everywhere and split at all places in $S$. There is one
place of $L_d$ over each of $3$ and $53$, with residue degrees one
and two respectively. Both are totally ramified in $L_{d+1}/L_d$.
The actual vertical ideal norms, extensions and automorphisms satisfy
\[
 \begin{gathered}
 N_d^L:C_{d+1}\xrightarrow{\sim}C_d,\qquad
 C_d\simeq C_9^2,\qquad \sigma_d|_{C_{d+1}}=1,\\
 m_2=1,\quad m_d=N_2^L\cdots N_{d-1}^L,\qquad
 m_dN_d^L=m_{d+1},\quad m_{d+1}j_d^L=3m_d,\\
 \varprojlim_{d,N_d^L}C_d\xrightarrow[\mathrm{pr}_2]{\sim}C_2,
 \qquad j_d^L=3(N_d^L)^{-1}.
 \end{gathered}
 \tag{E123.2}
\]
The inverse in the last formula is the inverse of the displayed
isomorphism. In the actual E.121 bases on $A_2=C_9\oplus C_3^3$
and $C_2=C_9u\oplus C_9a$, every horizontal norm and extension,
transported by $n_d$ and $m_d$, is
\[
 N_h=\begin{pmatrix}1&0&6&6\\0&3&6&3\end{pmatrix},
 \qquad
 j_h=\begin{pmatrix}0&3\\2&0\\2&0\\0&1\end{pmatrix}.
 \tag{E123.3}
\]
The Artin map for $B_d/L_d$, in these coordinates, is
$(r,s)\mapsto s\bmod3$, with the positive generator
$h(\xi)=\xi^4$.

For every $d\ge2$, all later one-step unit norm quotients are evaluated:
\[
 \begin{gathered}
 V_d/N_{L_{d+1}/L_d}V_{d+1}\simeq C_3^3,\qquad
 \frac{V_d\cap N_{L_{d+1}/L_d}L_{d+1}^{\times}}
      {N_{L_{d+1}/L_d}V_{d+1}}\simeq C_3^2,\\
 V_d/N_{B_d/L_d}U_d\simeq C_3,\qquad
 U_d=\mathcal O_{B_d,S}^{\times}.
 \end{gathered}
 \tag{E123.4}
\]
Their forward maps and inverses are the actual Hilbert--90 divisor
maps in the proof. The kernels after completion for multiplication by
$3^n$ are, respectively,
\[
 \ker(\widehat V_{d+1,3}\to\widehat V_{d,3})
       \simeq\mu_3\oplus\mathbb Z_3^{\,2\cdot3^d},
 \qquad
 \ker(\widehat U_{d,3}\to\widehat V_{d,3})
       \simeq\mu_3\oplus\mathbb Z_3^{\,2\cdot3^d+4}.
 \tag{E123.5}
\]
The torsion factors are literal roots of unity. The proof specifies
the chosen basis maps, inverse maps and their ambiguity.

Here is the entire coherent divisor image for this positional factor.
Define
\[
 \begin{gathered}
 \mathcal P=\{p\ne3,53:v_3(\operatorname{ord}_{\mathbb F_p^\times}(53))=1\},\\
 \mathcal R(\mathcal P)=
 \{(v_p)\in\textstyle\prod_{p\in\mathcal P}\mathbb Z_3:
       v_p\equiv0\pmod{3^n}\text{ for almost all }p,\text{ for every }n\}.
 \end{gathered}
 \tag{E123.6}
\]
Each $p\in\mathcal P$ has exactly one branch
$(\mathfrak p_d)_{d\ge2}$ of residue degree one that survives the
completed ideal-norm limit. Put $c_p=m_d[\mathfrak p_d]\in C_2$.
The class is independent of $d$. The sum below is finite modulo nine:
\[
 c:\mathcal R(\mathcal P)\longrightarrow C_2,\qquad
 c(v)=\sum_pc_pv_p,\qquad
 \mathcal D=\ker c.
 \tag{E123.7}
\]
The map $c$ is onto. Concrete pivot primes are $31$ and $103$.
At level two their residues are
\[
 \begin{array}{c|cc|cc}
 p&\beta_2\bmod\mathfrak p_2&\beta_1\bmod\mathfrak p_1&
 p_L(c_p)\in C_9\oplus C_3&c_p\bmod3\text{ in }(u,a)\\ \hline
 31&27&29&(4,1)&(0,1)\\
 103&80&90&(4,0)&(1,1)
 \end{array}
 \qquad p_L=\begin{pmatrix}6&1\\2&1\end{pmatrix}.
 \tag{E123.8}
\]
Thus $(c_{31},c_{103})$ itself is an actual basis of $C_2$.
The literal HNF ideals, seven full principal ideal equalities,
complete relation lattice and inverse coordinate maps are retained
in the certificate; the certification uses no upper-field BNF
class-group assumption.

Write $c_p=a_pc_{31}+b_pc_{103}$, choosing $a_p,b_p\in\{0,\ldots,8\}$,
and put $D_p=e_p-a_pe_{31}-b_pe_{103}$ for $p\ne31,103$.
Then the following maps are inverse continuous homomorphisms:
\[
 \begin{aligned}
 \mathbb Z_3^2\times\mathcal R(\mathcal P\setminus\{31,103\})
 &\longrightarrow\mathcal D,\\
 (A,B,t)&\longmapsto
       9Ae_{31}+9Be_{103}+\sum_{p\ne31,103}t_pD_p,\\
 v&\longmapsto
 \left(\frac{v_{31}+\sum_pa_pv_p}{9},
       \frac{v_{103}+\sum_pb_pv_p}{9},
       (v_p)_{p\ne31,103}\right).
 \end{aligned}
 \tag{E123.9}
\]
Both sums in the inverse omit $31,103$. All coefficients are bounded,
so these sums converge in $\mathbb Z_3$.

Put
$\mathcal X=\varprojlim_{d,N}\widehat{L_d^\times}_3$ and
$\mathcal V=\varprojlim_{d,N}\widehat V_{d,3}$.
There is an exact sequence of topological groups
\[
 0\longrightarrow\mathcal V\longrightarrow\mathcal X
       \xrightarrow{\operatorname{div}_{S^c}}\mathcal D
       \longrightarrow0.
 \tag{E123.10}
\]
For every basis vector in (E123.9), choose an actual coherent lift
$A_{31},A_{103},A_p$ with the indicated divisor.
Then
\[
 \begin{aligned}
 s(v)&=A_{31}^{A}A_{103}^{B}\prod_{p\ne31,103}A_p^{\,t_p},\\
 \mathcal V\times\mathcal D&\xrightarrow{\sim}\mathcal X,
       &(y,v)&\longmapsto y\,s(v),\\
 x&\longmapsto
       \bigl(x/s(\operatorname{div}_{S^c}x),
             \operatorname{div}_{S^c}x\bigr)
 \end{aligned}
 \tag{E123.11}
\]
are actual maps and inverse maps. Every infinite product converges
in each field-level power-completion topology. The compact group
$\mathcal V$ is torsion-free and is isomorphic to
$\prod_{i\ge1}\mathbb Z_3$ by the explicit character construction in
the proof. Its actual norm to $\widehat{\mathbb Q^\times}_3$ is retained,
rather than replaced by an abstract product coordinate.

More precisely, let $\mathfrak q_3,\mathfrak q_{53}$ be the two
$S$-primes of $L_2$, and define continuous homomorphisms
\[
 \begin{gathered}
 \nu_3(y)=v_{\mathfrak q_3}(y_2),\qquad
 \nu_{53}(y)=2v_{\mathfrak q_{53}}(y_2),\\
 \rho(v)=v_{\mathfrak q_3}(s(v)_2),\qquad
 \eta(v)=2v_{\mathfrak q_{53}}(s(v)_2),\\
 N_{\mathbb Q}(y\,s(v))
   =3^{\,\nu_3(y)+\rho(v)}53^{\,\nu_{53}(y)+\eta(v)}
       \prod_{p\in\mathcal P}p^{v_p}.
 \end{gathered}
 \tag{E123.12}
\]
The family $(\beta_d)_d$ is an actual element of $\mathcal V$ with
$(\nu_3,\nu_{53})=(0,2)$. Equation (E123.12), together with the
unchanged norms of the other positional factors, is the literal
total-norm map in E.46. In particular $\mathcal D$ embeds as a closed
family of actual $H$-section classes on every first-coordinate line
$3^k\kappa_{53}$ containing this factor: set this factor to $s(v)^3$
and all other factors to one. Its scaled second toric coordinate is
\[
       \rho(v)\kappa_3+\eta(v)\kappa_{53}
                 +\sum_{p\in\mathcal P}v_p\kappa_p .
 \tag{E123.13}
\]
The inverse on this family is the outside-$S$ divisor divided by
three. The proof keeps the factor map, the division by three and
the original total-norm cube condition explicit. The other positional
factors and the further intrinsic quotient maps in E.23 retain their
previous exact definitions.
\end{statement}

\begin{proof}
\emph{The fields and actual class-field maps.}
The degree of $L_d/K$ is $3^d$, by the valuation-one Eisenstein
equation at the prime over $53$. The same argument over
$\mathbb Q(\xi)$ gives $[B_d:L_d]=3$. For $d\ge2$,
$B_d=B_2L_d$, so E.121's actual unramified, $S$-split extension
$B_2/L_2$ base-changes to the stated extension $B_d/L_d$.
The automorphism $h$ cycles the three primes in each of E.122's
two triples. Its quotient therefore has one prime in each case;
the residue degrees and vertical total ramification follow from
the literal norm and restriction maps on these primes. This also
proves the asserted absence of other ramification.

Let $H_D^S$ denote the maximal unramified abelian three-extension
of $D$ in which every $S$-prime splits. The actual inclusions are
\[
 B_d\subset H_{L_d}^S\subset H_{B_d}^S,\qquad
 H_{L_d}^S=\text{the maximal subfield of }H_{B_d}^S
                        \text{ abelian over }L_d .
 \tag{E123.14}
\]
Indeed unramified and $S$-split extensions remain so under base
change. Conversely every abelian intermediate field over $L_d$
in the right side is unramified and $S$-split over $L_d$,
because $B_d/L_d$ has those properties.
For an extension of a group $A$ by a cyclic group generated by $h$,
with $A$ abelian, the commutator subgroup is $(h-1)A$:
choose one lift of $h$ and commute it with $A$; commutators of
its powers give no additional subgroup. Artin reciprocity and
(E123.14) thus give the exact sequence
\[
 0\longrightarrow (A_d)_h\xrightarrow{N_h}C_d
       \xrightarrow{\operatorname{Art}_{B_d/L_d}}C_3
       \longrightarrow0.
 \tag{E123.15}
\]
The arrow on the left is the actual ideal norm. Its agreement
with restriction of Frobenius is Milne V Proposition3.3, including
its proof. The class-field identification uses V Theorem3.6
and Example3.9 \cite{E_MilneCFT}.

E.122 transports the entire $h$-action to the matrix $H$ of E.121.
Its coinvariants have order $27$, so (E123.15) gives $|C_d|=81$
for all $d\ge2$. Moreover
$H_{L_d}^S\cap L_{d+1}=L_d$: any nontrivial intersection would be
the cyclic cubic $L_{d+1}/L_d$, which is ramified at both
$S$-places. For completeness, the actual norm surjectivity can
be read directly from this intersection. For a prescribed
Frobenius in $H_{L_d}^S/L_d$, choose its extension acting identically
on $L_{d+1}$ in the compositum, then choose a prime with that
Frobenius. A prime above it has relative residue degree one,
so its ideal norm has the prescribed class.
This is the exact norm/Frobenius square of Milne V3.3 and the
Chebotarev theorem VIII7.4. Equal orders now make every $N_d^L$
an isomorphism. The identity $N_d^L\sigma_d=N_d^L$ proves that
$\sigma_d$ acts identically; norm followed by extension is
the trace $1+\sigma_d+\sigma_d^2=3$. This proves (E123.2).

Norms commute in the field tower. Norm and extension in the
horizontal-vertical square commute because
$B_d\cap L_{d+1}=L_d$, the first extension being $S$-split
and the second totally ramified there, and their compositum is
$B_{d+1}$. On an ideal this equality is the equality of the
residue-degree weighted prime sums. Thus E.121's four complete
class maps transport to (E123.3). In particular their Artin,
kernel and image maps are retained; no inverse of a noninjective
norm is used.

There is also a useful literal local comparison. Write
$53=-(1-54)$ and define
$a_1=-\exp(\log(1-54)/9)\in\mathbb Q_3$.
The logarithm has valuation three, so $a_1^9=53$ and
$v_3(a_1+1)=1$. At every labeled completion of $B_2$ over three,
$(\beta_2/a_1)^9=\zeta_3$, hence $\beta_2/a_1$ is a primitive
twenty-seventh root of unity. The local degrees give
\[
            (B_2)_v=(L_2)_w=\mathbb Q_3(\mu_{27}).
 \tag{E123.16}
\]
This is a map of the actual completed fields. Globally
$\mu(L_d)\{3\}=\mu_3$: otherwise $\xi\in L_d$, contradicting
$[B_d:L_d]=3$. E.122 gives $\mu(B_d)\{3\}=\mu_9$.

\emph{An actual unit character, and every later unit quotient.}
Use the literal root $v_{20}$ of E.122 in
$\mathbb Q[b,z]/(b^9-53z^3,z^6+z^3+1)$, and put
\[
 w=v_{20}(b,z)v_{20}(b,z^4)v_{20}(b,z^7)\in L_2^\times.
 \tag{E123.17}
\]
The rational-ring certificate rechecks its full cube identity,
each horizontal substitution, and its complete norm. The output
gives every coefficient of the resulting polynomial of degree
at most seventeen in $b=\beta_2$. A separate no-input GP program
checks that polynomial by exact resultants. The decisive values are
\[
 \begin{gathered}
 N_{L_2/K}(w)=\frac{3+6\zeta_3}{53},\qquad
 N_{L_2/\mathbb Q}(w)=\frac{27}{53^2},\\
 v_{\mathfrak q_{53}}(w)=-1,\qquad
 \overline{w\beta_2}=10+50\zeta_3\in\mathbb F_{53^2},\qquad
 (10+50\zeta_3)^{936}=\zeta_3^2 .
 \end{gathered}
 \tag{E123.18}
\]
Its coefficient denominator divides $53^2$. Since its absolute
norm has support only at $3,53$, the whole principal divisor
is $S$-supported: away from $S$ the element is integral, and the
positive norm valuations there sum to zero. Thus $w\in V_2$.
The positive local Artin convention of E.122 is minus the
displayed tame residue logarithm, so $\chi_{53}(w)=1$.
The product formula gives $\chi_3(w)=2$. This also agrees with
the two triple sums of the literal root's six characters.

These two characters remain $(2,1)$ on the included $w\in V_d$
for every $d\ge2$. At either unique local branch,
$N_{(L_d)_v/(L_2)_w}\beta_d=\beta_2$. The cup-product projection
formula therefore reads
\[
 \operatorname{Cor}\bigl([\beta_d]\cup\operatorname{Res}[w]\bigr)
       =[\beta_2]\cup[w],\qquad
 \operatorname{inv}_{(L_2)_w}\operatorname{Cor}
       =\operatorname{inv}_{(L_d)_v}.
 \tag{E123.19}
\]
This is the actual restriction/corestriction map with common
coefficient $\mu_3$, not multiplication of a character by the
extension degree. Milne III3.1 proves the required compatibility
by cup products; V5.2 gives the product of local reciprocity maps.
The local norm quotient has order three, so the ordinary field-norm
obstruction on $V_d$ is exactly the one-dimensional character
$\chi_{53}$, with $\chi_3=-\chi_{53}$.
The cyclic Hasse norm theorem VIII3.1, with its injectivity proof,
identifies its kernel with $V_d\cap N L_{d+1}^{\times}$
\cite{E_MilneCFT}.

Here are the precise unit-defect maps, including the completion
step. For a cyclic cubic $E/D$ in this tower, write $\tau$ for the
positive generator and $V_E,V_D$ for its $S$-units. Given
$u=Nx\in V_D$, the norm of $\operatorname{div}_{S^c}(x)$ is zero.
The cyclic permutation-module calculation on every outside-$S$
prime orbit gives an outside-$S$ divisor $D_x$ with
\[
 \operatorname{div}_{S^c}(x)=(\tau-1)D_x,\qquad
 \nu(u)=\operatorname{pr}_3[D_x]\bmod j\operatorname{Cl}_S(D)\{3\}.
 \tag{E123.20}
\]
Changing $x$ with its norm fixed changes it by a Hilbert--90
coboundary, hence changes $D_x$ by a principal divisor.
Changing $D_x$ with $(\tau-1)D_x$ fixed adds an invariant divisor.
Outside $S$ the extension is unramified, so invariant divisors
are precisely extended base divisors. These are the exact
ambiguities in (E123.20).
Its image is
$\operatorname{Cl}_S(E)\{3\}^{\tau}/j\operatorname{Cl}_S(D)\{3\}$:
for an invariant class choose a representative $D$, choose $x$
with $(\tau-1)D=\operatorname{div}_{S^c}(x)$, and send it back
to $Nx$. Its kernel consists exactly of norms of $S$-units,
by dividing $x$ by the corresponding principal Hilbert--90
coboundary. The prime-to-three part of the invariant quotient
vanishes because norm followed by extension is the trace,
which is multiplication by three on invariant classes.
This proves both directions of the map, also for the
horizontal extension.

For the vertical extension, (E123.2) makes this invariant quotient
$C_2/3C_2=C_3^2$. For the horizontal extension all $S$-units
are field norms: at $S$ the extension is split, and elsewhere
a unit is a local norm in an unramified extension.
In the E.121 coordinates its invariant quotient is
\[
 A_d^h/j_hC_d\xrightarrow{\sim}\mathbb F_3,\qquad
 (3r,t,t,s)\longmapsto r-s.
 \tag{E123.21}
\]
An inverse is the class of $(3,0,0,0)$ followed by the literal
ideal construction in (E123.20).
To give coordinates on the full vertical quotient, first subtract
the unique power $w^e$, $e=\chi_{53}(v)\in\{0,1,2\}$.
Apply (E123.20) to $vw^{-e}$ and use either actual basis of
$C_2/3C_2$. Conversely choose divisors representing that basis,
choose the stated principal corrections $x_1,x_2$, and map
$(e,r,s)$ to the class of $w^e(Nx_1)^r(Nx_2)^s$.
All carries of three vanish because every included cube is an
actual unit norm. This proves (E123.4) and the inverse maps.

The groups of $S$-units are finitely generated. Completion for
multiplication by $3^n$ preserves their exact norm-image sequence:
diagonalize the integer map between the free parts, retain the
finite root groups, and complete each displayed cyclic factor.
The free ranks are $3^d+1$ for $V_d$ and $3^{d+1}+5$ for $U_d$.
The vertical norm on $\mu_3$ is cubing; the horizontal norm on
$\mu_9$ sends $\xi$ to $\xi^{1+4+7}=\zeta_3$.
This gives precisely (E123.5).
For a completed kernel $K$, its torsion is the stated literal
$\mu_3$; choose a $\mathbb Z_3$-basis of $K/\mu_3$ and lifts
$k_i\in K$. The isomorphism is
$(z,(a_i))\mapsto z\prod_i k_i^{a_i}$.
The inverse takes quotient coordinates and divides by their
lifted product. A change of basis acts by its actual invertible
matrix; changing lifts multiplies the first coordinate by the
corresponding product of roots of unity.

\emph{The two literal prime classes and an unconditional presentation.}
The lower field $L_1$ and its complete ideal and unit data are
certified in E.119. Its outside-$S$ ideal basis $G_1,G_2$ and local
Artin map give the following direct recipe. For a degree-one prime
$\mathfrak p_1$, compute the exact equality
$\mathfrak p_1=G_1^{e_1}G_2^{e_2}(x)$, including all $S$-valuations.
Its actual idelic coordinate is
$(-\chi_3(x),-\chi_{53}(x),e_1,e_2)$:
the sign follows by applying global reciprocity to $(x)$.
Project this vector through E.119's exact kernel lattice and
Smith map. For $31$ and $103$ the calculation gives
\[
 \begin{array}{c|c|c|c}
 p&(e_1,e_2)&(\chi_3(x),\chi_{53}(x))&
   (-\chi_3(x),-\chi_{53}(x),e_1,e_2)\\ \hline
 31&(5,0)&(1,1)&(-1,-1,5,0)\\
 103&(2,2)&(1,0)&(-1,0,2,2).
 \end{array}
 \tag{E123.22}
\]
Every rational coefficient of $x$, the whole lower HNF prime
ideal, and the full principal ideal equality are in the
replayable GP output. Applying the exact Smith map gives
(E123.8). Reduction of $p_L$ modulo three is invertible;
hence those two full classes are independent modulo three
and form a basis of the proved group $C_9^2$.

The residue condition identifies the actual surviving primes,
not just arbitrary primes of the same rational norm.
Both rational primes are $4$ modulo nine and $53$ is noncubic.
For $m=(p-1)/3$, the unique prime-to-three root is
$b_p=53^{\,3(27^{-1}\bmod m)}\bmod p$.
It satisfies $b_p^m=1$, $b_p^{27}=53^3$, and
$\zeta_3=b_p^9/53$. Thus $\beta_2=b_p$ specifies the desired
degree-one prime; its lower residue is $b_p^3$.
The output verifies these identities for the displayed numbers.

Here is a full literal certification of the upper presentation.
The proposed ordinary generators $I_1,I_2,I_3$ supply seven
checked ideal equalities: their powers $9,9,3$, the two $S$-prime
relations, and the two pivot-prime corrections.
No correctness assumption about the proposed BNF group is used.
Those equalities give an actual homomorphism from
\[
 \mathbb Z^3/
 \begin{pmatrix}9&0&0&6&0\\0&9&0&3&0\\0&0&3&1&0\end{pmatrix}
       \mathbb Z^5
 \longrightarrow\operatorname{Cl}_S(L_2).
 \tag{E123.23}
\]
The source has order $81$. Both independently proved generating
classes $c_{31},c_{103}$ lie in its image by the checked pivot
equalities, so the map is an isomorphism onto the entire
three-primary subgroup. This proves the required certification
without the interrupted upper-field certification attempt.
The exact maps between generator coordinates and the prime basis are
\[
 P=\begin{pmatrix}7&7&0\\2&4&3\end{pmatrix},\qquad
 Q=\begin{pmatrix}5&4\\8&5\\1&0\end{pmatrix},\qquad
 PQ=I_2\pmod9 .
 \tag{E123.24}
\]
The inverse composition differs from the identity by the relation
lattice in (E123.23). The full Smith matrices and this equality
are checked over the integers. All HNF matrices are in the
literal integral basis
$1,\zeta_3,\beta_2,\zeta_3\beta_2,\ldots,\beta_2^8,
\zeta_3\beta_2^8$.
It is maximal: the valuation-one Eisenstein equation at $53$
is tame of degree nine over $K$; (E123.16) proves maximality at
three, and the polynomial discriminant has no other prime support.
The discriminant of the stated basis is $-3^{45}53^{16}$,
as checked independently. The candidate object is used only to
propose exact ideal witnesses, each of which is verified.
A further standalone program reads no file and performs no BNF
operation: it reconstructs the two number fields from the literal
polynomials and bases, obtains empty \texttt{nfcertify} obstruction
lists, and rechecks all nine lower and upper principal ideal
equalities, both positive Artin columns, the chosen prime residues,
and the full forward and inverse relation maps.

For any other prime in $\mathcal P$, its three-primary coordinate
$c_p$ can therefore be evaluated by exact ideal relations in this
certified finite presentation. One unconditional terminating
procedure enumerates positive integers $m\equiv1\pmod9$, the $81$
prime-basis classes, integral powers of the two $S$-primes, and
nonzero field elements, until the full ideal equality for
$\mathfrak p_2^{\,m}$ is found. Some such $m$ annihilates the
prime-to-three component of $[\mathfrak p_2]$, by the Chinese
remainder theorem. On $C_2=C_9^2$ multiplication by $m$ is the
identity, so the recovered coordinate is exactly $c_p$ and is
unique. This is the explicit projection to the three-primary
subgroup; it does not require a value for the remaining class-group
part. A principal-ideal routine may propose the equality directly;
it is accepted only after full ideal verification. No missing
coordinate is assigned from the genus projection.

\emph{Surviving branches, completion and coherent lifts.}
At a prime away from $3,53$, Kummer roots identify prime branches
with Frobenius orbits on the compatible $3$-power roots of
$53\zeta_3$. On a fixed cyclotomic component the action is affine
on $\mathbb Z_3$. Its only bounded orbits are fixed points:
if a positive power fixes $x$, then
$(q^n-1)(x-x_0)=0$ in the torsion-free ring $\mathbb Z_3$.
The required divisibility is uniform on each finite node away
from the fixed point. After a finite refinement such a node is
a residue ball on which $v_3(x-x_0)=t$ is constant.
Replace the residue Frobenius by its square when necessary, so
$q\equiv1\pmod3$, and put $s=v_3(q-1)$.
The exact formula
$\phi^r(x)-x=(q^r-1)(x-x_0)$ and the elementary lifting-the-exponent
identity give orbit three-part
$3^{\max(0,m-s-t)}$ modulo $3^m$, uniformly on that ball.
If $x_0\notin\mathbb Z_3$, the valuation is constant on every
integral ball after the same refinement.
Thus the norm from sufficiently deep descendants of any fixed
nonfixed node has an arbitrarily large common three-power weight.
Its coherent integral mass is divisible by every power of three
and is zero. In the nested nodes containing $x_0$, all other
descendants vanish by this argument; precisely the fixed descendant
retains its mass.
For a split rational prime the root of unity component of $53$
is canceled by exactly one of $\zeta_3,\zeta_3^2$ precisely when
$v_3(\operatorname{ord}_p53)=1$; it then gives the unique fixed
degree-one branch. For an inert prime, the residue of $53$
has prime-to-three order while the factor $\zeta_3$ has order
three, so neither component has a fixed branch.
This proves (E123.6) and the stated unique branch.
Norms of its prime ideals have residue weight one, so their
classes agree under the actual maps $m_d$.

For clarity, the completed principal-divisor exactness is also
literal. At a fixed field level let $P$ be the group of principal
outside-$S$ divisors. It is a subgroup of a free abelian divisor
group and is free. Thus
$0\to V_d\to L_d^\times\to P\to0$ splits as abelian groups.
The quotient of the divisor group by $P$ is the finite
$S$-class group. Its three-primary exponent is nine for $d\ge2$.
If $x\in P\cap3^n\operatorname{Div}_{S^c}$, write $x=3^ny$.
Then $[y]$ is three-primary and killed by nine, so
$9y\in P$ and $x\in3^{n-2}P$ for $n\ge2$.
These two cofinal filtrations show that the completed image of
$P$ is exactly the kernel of the finite map to $C_d$.
The completed divisor group is the restricted sequence group,
not the full product with its product topology.

For $v\in\mathcal R(\mathcal P)$, the resulting divisor at every
level has class $c(v)$. Hence $c(v)=0$ is necessary.
It is sufficient as well. At each level the fibre of completed
elements with that divisor is a nonempty compact torsor under
$\widehat V_{d,3}$. Any finite chain of norm equations is solved
by choosing an element at the highest level and taking its
actual norms down the chain. The finite intersection property
in the product of these compact fibres gives a compatible
element at every level. This proves (E123.10); it does not
assume that unit norms are onto.
The explicit prime basis proves the surjectivity of $c$ and
both inverse formulas in (E123.9).

Choose coherent lifts of the divisor basis as in (E123.11).
Modulo $3^n$ at any fixed field level, all but finitely many
exponents $t_p$ vanish, so the infinite product converges.
This also proves continuity, and the displayed inverse is
continuous by division and the divisor map.
If chosen lifts are replaced by $A_i y_i$, $y_i\in\mathcal V$,
then $s'(v)=s(v)\prod_i y_i^{a_i(v)}$.
The corresponding coordinate change is
$(y,v)\mapsto(y\prod_i y_i^{-a_i(v)},v)$.
This identifies every coordinate change.

The compact metrizable group $\mathcal V$ is torsion-free.
A torsion coherent family has every component in the literal
$\mu_3$, and every further vertical norm is cubing, so every
fixed component is one. Its Pontryagin dual is therefore a
countable divisible three-primary group.
Choose a basis $\chi_{i,1}$ of its three-torsion vector space,
and recursively choose $3\chi_{i,n+1}=\chi_{i,n}$.
The map to a product of copies of $\mathbb Z_3$ is
\[
       \Phi(y)_i\bmod3^n=3^n\chi_{i,n}(y)
                    \quad\text{in }\mathbb Z/3^n.
 \tag{E123.25}
\]
The characters on the right have values in
$3^{-n}\mathbb Z/\mathbb Z$. Their independence and generation
follow by successively multiplying a finite-order character
by its highest possible power of three and subtracting its
basis lift. Thus the inverse is the unique point satisfying
all these finite character equations; compactness gives
existence. This proves the topological inverse, not merely
an equality of ranks.

There are countably infinitely many factors. Apply Kato
Theorem1.1(2) to $F=L_d$, $p=3$, and $a_1=53$:
the field has one place over three, and the actual extension
$L_d(\mu_{3^\infty},53^{1/3^\infty})$ contains all later $L_e$.
The theorem uses the coherent norm image as defined on its
first page and proves finite index in the completed group
of three-units. Restrict that actual coherent system to the
fields $L_e$. Its image in $V_d$ has rank $3^d$.
These unbounded ranks exclude a finite number of factors in
(E123.25) \cite{E_KatoUniversalNorms}.
Only this proved finite-index assertion is used; no index
or unproved numerical unit basis is assigned.

Finally, the rational norm of a prime ideal in the unique
surviving branch is $p$, and the two $S$-prime residue degrees
are one and two. This proves every exponent in (E123.12).
The family $\beta_d$ is norm-compatible because the constant
term of its odd-degree radical equation gives
$N_{L_{d+1}/L_d}\beta_{d+1}=\beta_d$.
For a first coordinate $3^k\kappa_{53}$, the positional
factor under discussion is the actual $j=1$ factor in
E.46/E.111, with root coordinate sent to $\beta_d$ along its
cofinal subsystem; root cubing is its field norm.
Thus, if $\mathcal X_{\ne1}$ denotes the other actual factors,
the full $H$-fibre with this factor expanded is exactly
\[
 \left\{(z,y,v)\in\mathcal X_{\ne1}\times\mathcal V\times\mathcal D:
 N(z)\,3^{\nu_3(y)+\rho(v)}53^{\nu_{53}(y)+\eta(v)}
               \prod_pp^{v_p}\in
                  (\widehat{\mathbb Q^\times}_3)^3\right\}.
 \tag{E123.26}
\]
This is an equality through E.46's cocycle--norm maps, with
their inverse, and retains the actual norms of the other
factors. Substitution $z=1$, $y=1$, $v\mapsto3v$
gives (E123.13). Its inverse is the outside-$S$ divisor divided
by three. In the product coordinates (E123.11) its image has
unit coordinate one and divisor in the closed subgroup
$3\mathcal D$, which proves closedness.
The subgroup $3\mathcal D$ need not be replaced by
$\mathcal D\cap3\mathcal R(\mathcal P)$; the displayed map
uses precisely the former. The computation evaluates this
positional factor, its coherent divisor image and the actual
section subfamily. The remaining valuation functionals and
other positional factors in (E123.26) are explicitly retained
as maps, and no additional intrinsic lift is inferred.

The manuscript anchor is the same Fermat-tower and arithmetic
splitting discussion, \S51.16, folios733v--736v, scans172--178,
used in E.23 and E.46. The code records all stopped attempts
separately, including the upper flag-one certification timeout;
the final ideal certification above uses the independent actual
basis (E123.8) and literal ideal equalities.
\end{proof}

